% Source: MQ34b_v1.html
\documentclass[11pt]{article}
\usepackage[utf8]{inputenc}
\usepackage[T1]{fontenc}
\usepackage[margin=1in]{geometry}
\usepackage{amsmath}
\usepackage{amssymb}
\usepackage{siunitx}
\usepackage{enumitem}
\usepackage{parskip}

\title{PS160 -- MQ34b\_v1.html}
\date{}

\begin{document}
\maketitle

\section*{Questions}

\begin{enumerate}[leftmargin=*]
  \item A lens is made of a material of index 1.9 with R1 = 10 cm and R2 = -12 cm. If the object is at 25 cm, what is the image position (in cm)? \hfill \textit{(1 pts, Numeric)}

  \item A coin is embedded in a crystal ball of radius 39 cm and index of refraction 1.3, 12 cm away from the surface. Calculate the image distance in centimeters. (Remember that the image distance can be either positive or negative!) \hfill \textit{(1 pts, Numeric)}

  \item The figure (not to scale) shows a tiny tree and a system of two converging lenses separated by a distance \emph{L}= 18 cm. The tree is 12 cm from lens 1. Find the \emph{x-}position in cm of the final image if \emph{f}\textsubscript{1} = 7 cm and \emph{f}\textsubscript{2} = 13 cm.

[Image: two\_lenses\_01.png] \hfill \textit{(1 pts, Numeric)}

  \item A man wears glasses with power -4.32 diopters so he can see distant objects. What is his far point? Answer in centimeters. (Give a positive number in cm for the distance in front of his glasses.) \hfill \textit{(1 pts, Numeric)}

  \item A simple magnifier has a focal length of 9.8 cm. Calculate the magnification. \hfill \textit{(1 pts, Numeric)}

\end{enumerate}

\end{document}